\section{Introduction}
\label{sec:introduction}

\subsection{Why a Second CVD Metric Exists}

Centroidal Voronoi Districts (\cvd{}, T.10~\citep{dellaLibera2026cvd}) are a
seed-and-grow split primitive for BISECT. At each bisection node, the algorithm
places two seeds, assigns every tract to its nearest seed, moves each seed to
the center of the tract set it owns, and repeats. The important choice is what
``nearest'' means.

Graph-distance CVD answers nearest by adjacency hops. That is often a good
redistricting default because contiguity is native to the graph: a tract one
edge away is close in the bisection graph, and the medoid of a district is an
actual tract. Geographic CVD answers nearest by projected coordinate distance:
each tract has a longitude/latitude centroid, BISECT projects those centroids
with its Albers transform, and Euclidean distance in meters decides ownership.

The two metrics see different worlds. A water body, mountain corridor, or long
string of small urban tracts can make graph distance and map distance disagree.
When they disagree, the bisection seed should not be chosen by accident; it
should be chosen by an explicit metric that an auditor can reproduce.

\subsection{Show the Difference}

The following schematic shows the difference between saying ``two tracts are
close'' in graph space and showing why geographic CVD may choose a different
split. The letters are tracts, not districts. The gap is water or empty land:
there is no adjacency edge across it, but the map distance across the gap is
short.

\begin{center}
\begin{tabular}{c@{\qquad}c}
\textbf{Map view} & \textbf{Graph-hop view} \\
\begin{tabular}{|c|c|c|c|c|c|}
\hline
A & A & A &   & B & B \\ \hline
A & A &   &   & B & B \\ \hline
A & A & A &   & B & B \\ \hline
\end{tabular}
&
\begin{tabular}{c}
\texttt{A--A--A--A--A--B--B--B--B--B}\\
\texttt{        land corridor only}
\end{tabular}
\end{tabular}
\end{center}

Graph-distance CVD grows through the available corridor. Geographic CVD sees
the projected coordinates and may assign the right shore to the right seed
even if the graph path between shores is long. This does not replace contiguity
checking; it changes the proximity signal used before BISECT's local rebalance
and final validation.

\subsection{Implemented Scope}

This paper documents the implemented BISECT surface:

\begin{itemize}
\item Graph-distance CVD uses BFS hop count and graph-distance medoids.
\item Geographic CVD uses projected tract centroids, Euclidean assignment, and
  population-weighted coordinate updates snapped to the nearest tract.
\item \texttt{"CVD\_GEO\_INIT\_"} domain-separates geographic CVD node seeds
  from graph-distance CVD seeds.
\item Missing centroid data is a hard configuration error for geographic CVD,
  not a silent fallback to graph-distance CVD.
\end{itemize}

The current implementation does not yet provide state-specific projection
dispatch, archived NC/FL/WA comparison packages, or ratio-aware CVD target
populations for odd-seat bisection nodes. Those items are the evidence ceiling
for this paper.

\subsection{Relationship to T.10}
\label{sec:relationship-b22}

T.10 is the canonical CVD paper: it explains proximity packing, the
graph-distance implementation, and the shared recursive bisection contract.
This T.11 note is the geographic-metric companion. It is narrower by design:
where T.10 asks ``what is CVD?'', this paper asks ``what changes when CVD uses
projected coordinate distance, and how does an auditor reproduce that choice?''

\subsection{Relation to Prior Work}

The continuous CVT framework of Du, Faber, and Gunzburger~\citeyear{du1999}
applies Lloyd's algorithm in Euclidean space. Geographic CVD is a discrete,
tract-level adaptation: the centroid is the population-weighted mean of tract
coordinates, and the next seed is the real tract nearest to that mean.
Snyder~\citeyear{snyder1987} gives the Albers projection formulas used by the
platform's spherical transform; the present paper treats projection metadata as
part of the audit package rather than as an incidental implementation detail.
